\documentclass[10pt,a4paper]{amsart}
\usepackage{amsmath,amssymb,amsthm,graphicx}
% Retained preamble for the Stacks Project chapter morphisms.tex (Morphisms of Schemes).
% The chapter file pinned-original-source.tex carries only \input{preamble} on its first line, so no
% preamble text lives inside the chapter itself.  The Stacks Project uses one single shared preamble for
% every chapter; the snapshot stages it verbatim in sources/stacks-chow/preamble.tex.  The block below is
% copied mechanically, line for line, from that file's lines 14-18 (the Xy-pic setup) and lines 154-217
% (the native theorem declarations together with the mathematical macros).  Nothing was added, removed or
% redefined.  The wrapper class and the amsmath/amssymb/amsthm/graphicx base packages are supplied by the
% assembler, so the class/\externaldocument part of the shared preamble is not repeated here.
% For commutative diagrams we use Xy-pic
\usepackage[all]{xy}

% We use 2cell for 2-commutative diagrams.
\xyoption{2cell}
% Theorem environments.
%
\theoremstyle{plain}
\newtheorem{theorem}[subsection]{Theorem}
\newtheorem{proposition}[subsection]{Proposition}
\newtheorem{lemma}[subsection]{Lemma}

\theoremstyle{definition}
\newtheorem{definition}[subsection]{Definition}
\newtheorem{example}[subsection]{Example}
\newtheorem{exercise}[subsection]{Exercise}
\newtheorem{situation}[subsection]{Situation}

\theoremstyle{remark}
\newtheorem{remark}[subsection]{Remark}
\newtheorem{remarks}[subsection]{Remarks}

\numberwithin{equation}{subsection}

% Macros
%
\def\lim{\mathop{\mathrm{lim}}\nolimits}
\def\colim{\mathop{\mathrm{colim}}\nolimits}
\def\Spec{\mathop{\mathrm{Spec}}}
\def\Hom{\mathop{\mathrm{Hom}}\nolimits}
\def\Ext{\mathop{\mathrm{Ext}}\nolimits}
\def\SheafHom{\mathop{\mathcal{H}\!\mathit{om}}\nolimits}
\def\SheafExt{\mathop{\mathcal{E}\!\mathit{xt}}\nolimits}
\def\Sch{\mathit{Sch}}
\def\Mor{\mathop{\mathrm{Mor}}\nolimits}
\def\Ob{\mathop{\mathrm{Ob}}\nolimits}
\def\Sh{\mathop{\mathit{Sh}}\nolimits}
\def\NL{\mathop{N\!L}\nolimits}
\def\CH{\mathop{\mathrm{CH}}\nolimits}
\def\proetale{{pro\text{-}\acute{e}tale}}
\def\etale{{\acute{e}tale}}
\def\QCoh{\mathit{QCoh}}
\def\Ker{\mathop{\mathrm{Ker}}}
\def\Im{\mathop{\mathrm{Im}}}
\def\Coker{\mathop{\mathrm{Coker}}}
\def\Coim{\mathop{\mathrm{Coim}}}

% Boxtimes
%
\DeclareMathSymbol{\boxtimes}{\mathbin}{AMSa}{"02}

%
% Macros for moduli stacks/spaces
%
\def\QCohstack{\mathcal{QC}\!\mathit{oh}}
\def\Cohstack{\mathcal{C}\!\mathit{oh}}
\def\Spacesstack{\mathcal{S}\!\mathit{paces}}
\def\Quotfunctor{\mathrm{Quot}}
\def\Hilbfunctor{\mathrm{Hilb}}
\def\Curvesstack{\mathcal{C}\!\mathit{urves}}
\def\Polarizedstack{\mathcal{P}\!\mathit{olarized}}
\def\Complexesstack{\mathcal{C}\!\mathit{omplexes}}
% \Pic is the operator that assigns to X its picard group, usage \Pic(X)
% \Picardstack_{X/B} denotes the Picard stack of X over B
% \Picardfunctor_{X/B} denotes the Picard functor of X over B
\def\Pic{\mathop{\mathrm{Pic}}\nolimits}
\def\Picardstack{\mathcal{P}\!\mathit{ic}}
\def\Picardfunctor{\mathrm{Pic}}
\def\Deformationcategory{\mathcal{D}\!\mathit{ef}}

% --- appended 2026-10-12 (Chief): the project own definitions of the reference/slogan/history
% display environments, copied mechanically from the shared Stacks preamble lines 8-12
% (\usepackage{verbatim} plus the three \newenvironment lines, which the chapter preamble keeps
% verbatim from the project preamble). Nothing added, removed or redefined: these environments
% carry no mathematical content and the project itself typesets them as comments. Without these
% three lines a counted statement whose own text uses them fails with Environment slogan undefined.
\usepackage{verbatim}
\newenvironment{reference}{\comment}{\endcomment}
\newenvironment{slogan}{\comment}{\endcomment}
\newenvironment{history}{\comment}{\endcomment}
\usepackage{cleveref}
\crefname{originallemma}{original lemma}{original lemmas}
\Crefname{originallemma}{Original lemma}{Original lemmas}
\makeatletter
\newlabel{lemma-finite-flat}{{02KB}{source}{}{sourcepointer-1}{}}
\newlabel{lemma-finite-flat@cref}{{[originallemma][0][]02KB}{[1][0][]source}{}{}{}}
\makeatother
\crefname{originallemma}{original lemma}{original lemmas}
\Crefname{originallemma}{Original lemma}{Original lemmas}
\makeatletter
\newlabel{lemma-finite-proper}{{01WN}{source}{}{sourcepointer-2}{}}
\newlabel{lemma-finite-proper@cref}{{[originallemma][0][]01WN}{[1][0][]source}{}{}{}}
\makeatother
\crefname{originallemma}{original lemma}{original lemmas}
\Crefname{originallemma}{Original lemma}{Original lemmas}
\makeatletter
\newlabel{lemma-fppf-open}{{01UA}{source}{}{sourcepointer-3}{}}
\newlabel{lemma-fppf-open@cref}{{[originallemma][0][]01UA}{[1][0][]source}{}{}{}}
\makeatother
\crefname{originallemma}{original lemma}{original lemmas}
\Crefname{originallemma}{Original lemma}{Original lemmas}
\makeatletter
\newlabel{lemma-massage-finite}{{03HW}{source}{}{sourcepointer-4}{}}
\newlabel{lemma-massage-finite@cref}{{[originallemma][0][]03HW}{[1][0][]source}{}{}{}}
\makeatother
\crefname{originallemma}{original lemma}{original lemmas}
\Crefname{originallemma}{Original lemma}{Original lemmas}
\makeatletter
\newlabel{topology-lemma-closed-image-nowhere-dense}{{03HQ}{source}{}{sourcepointer-5}{}}
\newlabel{topology-lemma-closed-image-nowhere-dense@cref}{{[originallemma][0][]03HQ}{[1][0][]source}{}{}{}}
\makeatother
\crefname{originallemma}{original lemma}{original lemmas}
\Crefname{originallemma}{Original lemma}{Original lemmas}
\makeatletter
\newlabel{topology-lemma-nowhere-dense}{{03HO}{source}{}{sourcepointer-6}{}}
\newlabel{topology-lemma-nowhere-dense@cref}{{[originallemma][0][]03HO}{[1][0][]source}{}{}{}}
\makeatother
\crefname{originallemma}{original lemma}{original lemmas}
\Crefname{originallemma}{Original lemma}{Original lemmas}
\makeatletter
\newlabel{topology-lemma-nowhere-dense-local}{{03HP}{source}{}{sourcepointer-7}{}}
\newlabel{topology-lemma-nowhere-dense-local@cref}{{[originallemma][0][]03HP}{[1][0][]source}{}{}{}}
\makeatother
\crefname{originallemma}{original lemma}{original lemmas}
\Crefname{originallemma}{Original lemma}{Original lemmas}
\makeatletter
\newlabel{topology-lemma-open-inverse-image-closed-nowhere-dense}{{03HR}{source}{}{sourcepointer-8}{}}
\newlabel{topology-lemma-open-inverse-image-closed-nowhere-dense@cref}{{[originallemma][0][]03HR}{[1][0][]source}{}{}{}}
\makeatother
\title{The image of a closed nowhere dense subset under a finite morphism is again closed and nowhere dense}
\author{The Stacks Project Authors}
\date{stacks-project:morphisms.tex:local-snapshot-20261012; GNU FDL 1.2 or later}
\begin{document}
\maketitle
\noindent\textit{Editorial note.} All delivered mathematical text is the author's own, copied from the pinned Stacks Project chapter file pinned-original-source.tex; no statement or proof step was rewritten, reordered or removed.\par\medskip
\noindent\textit{Editorial note.} Following the multi-claim convention for this source, source\_url is the source-level URL https://stacks.math.columbia.edu/ and the source label lemma-image-nowhere-dense-finite is carried in claim\_key and here rather than in a per-claim tag URL.\par\medskip
\noindent\textit{Editorial note.} The preamble is the chapter's shared Stacks Project preamble, retained as sources/stacks-morphisms/retained-preamble.tex: a line-for-line excerpt of the shared Stacks preamble staged with this snapshot. No macro was invented or redefined, and this card adds no preamble text of its own.\par\medskip
\noindent\textit{Editorial note.} Statement span 12499-12504 and proof span 12506-12570 come from the halfA pool card operations/pool-cards/stacks-morphisms-halfA.json and were re-read in the pinned chapter file before this card was written.\par\medskip
\noindent\textit{Editorial note.} The proof contains an xymatrix diagram with five arrows; it is copied verbatim and the retained chapter preamble loads Xy-pic.\par\medskip
\noindent\textit{Editorial note.} The proof's closing phrase 'and we win' is the author's and is preserved.\par\medskip
\noindent\textit{Editorial note.} 8 cross-references occur inside the counted spans; they are listed in external\_references as navigation pointers only.\par\medskip
\noindent\textit{Allowed standard prerequisites (editorial).}\begin{itemize}
\item A finite morphism is proper, hence closed, and the inverse image of a closed nowhere dense subset under an open morphism is nowhere dense. Cited at proof lines 12506-12508 and 12524-12529. Gives closedness of the image and lets the nowhere density be transported along the open morphisms appearing in the auxiliary diagram.
\item Nowhere density is local on the target, and a finite union of closed nowhere dense subsets is closed and nowhere dense. Cited at proof lines 12511-12513 and 12545-12548. Reduces the statement to an affine base and assembles the final case from finitely many closed pieces.
\item An open inverse image of a closed subset is closed, and the image of a closed nowhere dense subset under a closed immersion stays closed and nowhere dense. Cited at proof lines 12530-12533. Allows the auxiliary diagram to be replaced by the situation where the source is a finite union of closed pieces mapping isomorphically to the base.
\item The auxiliary diagram for a finite morphism over an affine base, with a finite locally free morphism to the target and to the source, cited from the same chapter. Cited at proof lines 12516-12518. Produces the open morphisms along which the problem is invariant.
\item A finite locally free morphism is flat, and a flat morphism locally of finite presentation is open, cited from the same chapter. Cited at proof lines 12521-12524. Shows that the morphisms of the auxiliary diagram are open.
\end{itemize}\medskip
\ifdefined\hypertarget\hypertarget{sourcepointer-1}{}\fi
\noindent\textit{Source reference.} Original lemma 02KB: focused-seeds/AGP/sources/stacks-snapshot-20261011/tags.txt, source lines 3366--3366. Same-chapter reference used in the retained author text at source line(s) 12527, 13544. It is supplied only as a navigation pointer so that the retained text compiles; the counted item delivers the author's own statement and the author's complete proof, and no part of the referenced item is copied in or substituted for the counted argument.\par\medskip
\ifdefined\hypertarget\hypertarget{sourcepointer-2}{}\fi
\noindent\textit{Source reference.} Original lemma 01WN: focused-seeds/AGP/sources/stacks-snapshot-20261011/tags.txt, source lines 2488--2488. Same-chapter reference used in the retained author text at source line(s) 11282, 12507, 13352. It is supplied only as a navigation pointer so that the retained text compiles; the counted item delivers the author's own statement and the author's complete proof, and no part of the referenced item is copied in or substituted for the counted argument.\par\medskip
\ifdefined\hypertarget\hypertarget{sourcepointer-3}{}\fi
\noindent\textit{Source reference.} Original lemma 01UA: focused-seeds/AGP/sources/stacks-snapshot-20261011/tags.txt, source lines 2403--2403. Same-chapter reference used in the retained author text at source line(s) 4492, 4621, 5693, 5890, 7179, 8281, 12527. It is supplied only as a navigation pointer so that the retained text compiles; the counted item delivers the author's own statement and the author's complete proof, and no part of the referenced item is copied in or substituted for the counted argument.\par\medskip
\ifdefined\hypertarget\hypertarget{sourcepointer-4}{}\fi
\noindent\textit{Source reference.} Original lemma 03HW: focused-seeds/AGP/sources/stacks-snapshot-20261011/tags.txt, source lines 4582--4582. Same-chapter reference used in the retained author text at source line(s) 12525. It is supplied only as a navigation pointer so that the retained text compiles; the counted item delivers the author's own statement and the author's complete proof, and no part of the referenced item is copied in or substituted for the counted argument.\par\medskip
\ifdefined\hypertarget\hypertarget{sourcepointer-5}{}\fi
\noindent\textit{Source reference.} Original lemma 03HQ: focused-seeds/AGP/sources/stacks-snapshot-20261011/tags.txt, source lines 4576--4576. Same-chapter reference used in the retained author text at source line(s) 12535. It is supplied only as a navigation pointer so that the retained text compiles; the counted item delivers the author's own statement and the author's complete proof, and no part of the referenced item is copied in or substituted for the counted argument.\par\medskip
\ifdefined\hypertarget\hypertarget{sourcepointer-6}{}\fi
\noindent\textit{Source reference.} Original lemma 03HO: focused-seeds/AGP/sources/stacks-snapshot-20261011/tags.txt, source lines 4574--4574. Same-chapter reference used in the retained author text at source line(s) 12547, 15267. It is supplied only as a navigation pointer so that the retained text compiles; the counted item delivers the author's own statement and the author's complete proof, and no part of the referenced item is copied in or substituted for the counted argument.\par\medskip
\ifdefined\hypertarget\hypertarget{sourcepointer-7}{}\fi
\noindent\textit{Source reference.} Original lemma 03HP: focused-seeds/AGP/sources/stacks-snapshot-20261011/tags.txt, source lines 4575--4575. Same-chapter reference used in the retained author text at source line(s) 12513. It is supplied only as a navigation pointer so that the retained text compiles; the counted item delivers the author's own statement and the author's complete proof, and no part of the referenced item is copied in or substituted for the counted argument.\par\medskip
\ifdefined\hypertarget\hypertarget{sourcepointer-8}{}\fi
\noindent\textit{Source reference.} Original lemma 03HR: focused-seeds/AGP/sources/stacks-snapshot-20261011/tags.txt, source lines 4577--4577. Same-chapter reference used in the retained author text at source line(s) 12529, 12533. It is supplied only as a navigation pointer so that the retained text compiles; the counted item delivers the author's own statement and the author's complete proof, and no part of the referenced item is copied in or substituted for the counted argument.\par\medskip
\begin{lemma}
\label{lemma-image-nowhere-dense-finite}
Let $f : Y \to X$ be a finite morphism of schemes.
Let $T \subset Y$ be a closed nowhere dense subset of $Y$.
Then $f(T) \subset X$ is a closed nowhere dense subset of $X$.
\end{lemma}

\begin{proof}
By Lemma \ref{lemma-finite-proper} we know that $f(T) \subset X$ is closed.
Let $X = \bigcup X_i$ be an affine covering.
Since $T$ is nowhere dense in $Y$, we see that also $T \cap f^{-1}(X_i)$
is nowhere dense in $f^{-1}(X_i)$. Hence if we can prove the theorem in the
affine case, then we see that $f(T) \cap X_i$ is nowhere dense.
This then implies that $T$ is nowhere dense in $X$ by
Topology, Lemma \ref{topology-lemma-nowhere-dense-local}.

\medskip\noindent
Assume $X$ is affine. Choose a diagram
$$
\xymatrix{
Z' \ar[rd] &
Y' \ar[l]^i \ar[d]^{f'} \ar[r]_a &
Y \ar[d]^f \\
 & X' \ar[r]^b & X
}
$$
as in Lemma \ref{lemma-massage-finite}. The morphisms $a$, $b$ are
open since they are finite locally free
(Lemmas \ref{lemma-finite-flat} and \ref{lemma-fppf-open}).
Hence $T' = a^{-1}(T)$ is nowhere dense, see
Topology, Lemma \ref{topology-lemma-open-inverse-image-closed-nowhere-dense}.
The morphism $b$ is surjective and open.
Hence, if we can prove $f'(T') = b^{-1}(f(T))$ is
nowhere dense, then $f(T)$ is nowhere dense, see
Topology, Lemma \ref{topology-lemma-open-inverse-image-closed-nowhere-dense}.
As $i$ is a closed immersion, by
Topology, Lemma \ref{topology-lemma-closed-image-nowhere-dense}
we see that $i(T') \subset Z'$ is closed and nowhere dense.
Thus we have reduced the problem to the case discussed
in the following paragraph.

\medskip\noindent
Assume that $Y = \bigcup_{i = 1, \ldots, n} Y_i$ is a finite union of
closed subsets, each mapping isomorphically to $X$. Consider
$T_i = Y_i \cap T$. If each of the $T_i$ is nowhere dense in $Y_i$,
then each $f(T_i)$ is nowhere dense in $X$ as $Y_i \to X$ is an isomorphism.
Hence $f(T) = f(T_i)$ is a finite union of nowhere dense closed
subsets of $X$ and we win, see
Topology, Lemma \ref{topology-lemma-nowhere-dense}.
Suppose not, say $T_1$ contains a nonempty open $V \subset Y_1$.
We are going to show this leads to a contradiction.
Consider $Y_2 \cap V \subset V$. This is either
a proper closed subset, or equal to $V$. In the first case we replace
$V$ by $V \setminus V \cap Y_2$, so $V \subset T_1$ is open in $Y_1$ and
does not meet $Y_2$. In the second case we have
$V \subset Y_1 \cap Y_2$ is open in both $Y_1$ and $Y_2$.
Repeat sequentially with $i = 3, \ldots, n$. The result is a disjoint
union decomposition
$$
\{1, \ldots, n\} = I_1 \amalg I_2, \quad 1 \in I_1
$$
and an open $V$ of $Y_1$ contained in $T_1$ such that $V \subset Y_i$
for $i \in I_1$ and $V \cap Y_i = \emptyset$ for $i \in I_2$. Set
$U = f(V)$. This is an open of $X$ since $f|_{Y_1} : Y_1 \to X$ is
an isomorphism. Then
$$
f^{-1}(U) = V\ \amalg\ \bigcup\nolimits_{i \in I_2} (Y_i \cap f^{-1}(U))
$$
As $\bigcup_{i \in I_2} Y_i$ is closed, this implies that
$V \subset f^{-1}(U)$ is open, hence $V \subset Y$ is open.
This contradicts the assumption that $T$ is nowhere dense in $Y$, as desired.
\end{proof}

\end{document}
